\section{Related Work}

Friedberg's 1957 paper constructed two recursively enumerable sets of
incomparable degrees and solved Post's problem \cite{friedberg1957}.
Muchnik independently obtained the same solution in 1956, with an expanded
publication appearing later \cite{muchnik1956,muchnik1958}.  Soare's
monograph remains a standard reference for recursively enumerable sets,
degrees, and priority methods \cite{soare1987}.  The formalization follows
the finite-injury presentation rather than a reformulation that avoids
priority machinery.

The Lean development builds on Mathlib's ordinary computability theory.
Carneiro formalized primitive recursive and partial recursive functions in
Lean, including universal partial recursive functions and the halting
problem \cite{carneiro2019}.  Mathlib's current
\texttt{Partrec} and \texttt{PartrecCode} libraries provide the closure
theorems and code infrastructure reused here
\cite{mathlibPartrec,mathlibPartrecCode}.  The present project differs by
adding a local relative oracle-program semantics with explicit finite uses,
which is tailored to priority reasoning.

Other proof assistants also contain computability developments.  The
Isabelle Archive of Formal Proofs includes a universal Turing machine
formalization with undecidability and computability-theoretic material
\cite{afpUTM}, building on Isabelle/HOL's mature theorem-proving
infrastructure \cite{nipkow2002}.  Coq has been used to formalize abstract
computability through Turing categories \cite{cockett2018}.  These works
address foundational computability models; the contribution here is a
direct mechanization of a finite-injury priority argument.

We are not aware of another publicly available Lean formalization of the
complete Friedberg--Muchnik finite-injury construction.  This claim should
be treated cautiously: a submission should include a more systematic search
of formalization archives and unpublished developments before making a
strong priority claim.

